% BrokerGuard v3 -- Springer LNCS, mục tiêu 15 trang kể cả tài liệu tham khảo.
% BẢN TIẾNG VIỆT ĐỂ ĐỌC (19/11/2026). Bản tiếng Anh viết sau, cùng cấu trúc, cùng bảng, hình và số.
% Mọi con số lấy từ ../../iccies2027/facts.md (mục Test-4, các mục bổ sung 19/11 và "Bổ sung cho bài v3").
% Overleaf: tải cả thư mục latex/ (zip), main = main.tex, trình biên dịch pdfLaTeX. Hình 1--3 vẽ bằng TikZ/pgfplots,
% Overleaf tự vẽ, không có file ảnh ngoài.
% Nếu pdfLaTeX báo lỗi font tiếng Việt (T5), đổi sang XeLaTeX và thay 3 dòng fontenc/inputenc/babel
% bằng: \usepackage{fontspec} \usepackage[vietnamese]{babel}
\documentclass[runningheads]{llncs}

\usepackage[T5]{fontenc}
\usepackage[utf8]{inputenc}
\usepackage[vietnamese]{babel}
\usepackage{graphicx}
\usepackage{booktabs}
\usepackage{multirow}
\usepackage{amsmath,amssymb}
\usepackage{xcolor}
\usepackage{url}
\usepackage{tabularx}    % bảng trải đúng bề ngang trang, cột chữ tự xuống dòng
\usepackage{makecell}    % tiêu đề cột hai dòng
\usepackage[section]{placeins} % bảng/hình không trôi sang mục sau
\usepackage{pgfplots}    % Hình 2 và 3
\pgfplotsset{compat=1.18}
\usetikzlibrary{patterns,arrows.meta}
% Bảng màu phân loại đã kiểm mù màu (xanh/cam ΔE CVD 24,7). Cam = \ours (thành phần mới), xanh = BG-Base,
% xám = baseline. Màu theo thực thể, giống nhau ở mọi hình.
\definecolor{bgblue}{HTML}{2A78D6}
\definecolor{bgorange}{HTML}{EB6834}
\usepackage{enumitem}
\setlist{noitemsep, topsep=2pt}

% Đặt bảng/hình gần chỗ nhắc tới, bớt khoảng trắng (không đổi lề, cỡ chữ hay giãn dòng của LNCS)
\renewcommand{\topfraction}{0.85}
\renewcommand{\textfraction}{0.1}
\renewcommand{\floatpagefraction}{0.75}
\setcounter{topnumber}{3}
\setcounter{totalnumber}{4}
\setlength{\textfloatsep}{8pt plus 2pt minus 2pt}
\setlength{\floatsep}{6pt plus 2pt minus 2pt}
\setlength{\intextsep}{6pt plus 2pt minus 2pt}
\makeatletter
\setlength{\@fptop}{0pt}
\setlength{\@fpsep}{10pt plus 2pt}
\setlength{\@fpbot}{0pt plus 1fil}
\g@addto@macro\normalsize{%
  \setlength{\abovedisplayskip}{4pt plus 2pt minus 1pt}%
  \setlength{\belowdisplayskip}{4pt plus 2pt minus 1pt}%
  \setlength{\abovedisplayshortskip}{2pt plus 1pt}%
  \setlength{\belowdisplayshortskip}{2pt plus 1pt}}
\makeatother

% Tên môi trường định lý và từ khóa bằng tiếng Việt (bản tiếng Anh: xóa khối này)
\addto\captionsvietnamese{%
  \renewcommand{\propositionname}{Mệnh đề}%
  \renewcommand{\definitionname}{Định nghĩa}%
  \renewcommand{\proofname}{Chứng minh}}
\renewcommand{\keywordname}{{\bfseries Từ khóa:}}

% Ghi chú cho tác giả; xóa trước khi nộp.
\newcommand{\todo}[1]{\textcolor{red}{[TODO: #1]}}
\newcommand{\code}[1]{\texttt{#1}}

% ---- TÊN GỌI (đổi một dòng ở đây là đổi cả bài; ánh xạ sang log ở ../README.md) -------------------
\newcommand{\ours}{\textsc{BrokerGuard}}     % v3 = V3P2: BG-Base + C1 + C2 + C3
\newcommand{\base}{BG-Base}                  % v1 = D2P2 (final_config_v3), luật chuỗi + Sandwich
\newcommand{\cRules}{BG-Rules}               % D2 = BG-Base không có lời nhắc
\newcommand{\cNone}{NoDef}                   % A0
\newcommand{\cSpot}{Spotlight}               % P1
\newcommand{\cSand}{Sandwich}                % P2
\newcommand{\cPG}{PromptGuard}               % G1 (Prompt Guard 2, 86M)
\newcommand{\cMel}{MELON$^*$}                % MEL (cài lại)
\newcommand{\cProg}{Progent$^*$}             % PRG (cài lại)
\newcommand{\cPhrase}{PhraseRules}           % R1
% Ba thành phần mới
\newcommand{\cOne}{\textsf{C1}}              % chính sách ý định P(u)
\newcommand{\cTwo}{\textsf{C2}}              % nhãn mức giá trị
\newcommand{\cThree}{\textsf{C3}}            % ghim tham số
\newcommand{\abl}[1]{\ours$-$\textsf{C#1}}  % bản bỏ một thành phần
% Luật giữ từ BG-Base
\newcommand{\rDest}{\textsf{R-Dest}}         % luật 2 của D2 (bị C1 thay)
\newcommand{\rFlow}{\textsf{R-Flow}}         % luật 3 của D2 (theo nhãn nguồn)
\newcommand{\rRedact}{\textsf{R-Redact}}     % luật 4 của D2
\newcommand{\rSig}{\textsf{R-Signal}}        % lớp A1 của D2
% Cơ chế tấn công (chỉ dùng trong bảng và hình)
\newcommand{\mER}{\textsf{ER}}               % M1
\newcommand{\mCR}{\textsf{CR}}               % M2
\newcommand{\mSL}{\textsf{SL}}               % M3
\newcommand{\mAL}{\textsf{AL}}               % M4
\newcommand{\mPT}{\textsf{PT}}               % M5
\newcommand{\mUS}{\textsf{US}}               % M6
% Tập dữ liệu
\newcommand{\dev}{Dev}                       % Dev v4.1 (168 họ)
\newcommand{\tOne}{Test-I}                   % "Test v3", nay thuộc Dev
\newcommand{\tTwo}{Test-II}                  % "Test-3", nay thuộc Dev
\newcommand{\tThree}{Test-III}               % "Test-4" (T5_*), tập kiểm định duy nhất của bài
% Ký hiệu
\newcommand{\sub}{\sqsubseteq}               % x ⊑ z: kappa(x) là chuỗi con của kappa(z)
\newcommand{\dest}{\mathit{dst}}
\newcommand{\data}{\mathit{data}}
\newcommand{\norm}{\nu}
\newcommand{\bPlain}{lành tính thường}       % benign_plain
\newcommand{\bAuth}{lành tính ủy quyền}       % benign_auth

% hyperref nạp cuối cùng: trích dẫn, tham chiếu và DOI/URL màu xanh, bấm được
\usepackage{hyperref}
\definecolor{linkblue}{HTML}{1F4E9A}
\hypersetup{colorlinks=true, linkcolor=linkblue, citecolor=linkblue, urlcolor=linkblue}
\AtBeginDocument{\def\doi#1{\href{https://doi.org/#1}{doi:#1}}}

\begin{document}

% Tiêu đề (tiếng Anh, là tiêu đề nộp). Nghĩa: BrokerGuard: cấp quyền theo ý định và kiểm soát luồng mức giá trị
% trước tấn công chèn lệnh gián tiếp cho tác tử LLM tiếng Việt.
\title{BrokerGuard: Intent-Aware Authorisation and Value-Level Flow Control Against Indirect Prompt Injection in
Vietnamese LLM Agents}
\titlerunning{BrokerGuard}

\author{\todo{Tác giả 1}\inst{1} \and \todo{Tác giả 2}\inst{1}}
\authorrunning{\todo{Tác giả} et al.}
\institute{\todo{Đơn vị, thành phố, quốc gia} \\ \email{\todo{email}}}

\maketitle

% LNCS: tóm tắt 150--250 từ (tính trên bản tiếng Anh). Không đưa giá trị p vào tóm tắt.
% Số: facts.md mục Test-4 (40, 6, 2 họ; cặp lệch 4/0; cùng chiều trên ba mô hình; ΔU −3,1; token và độ trễ v3 so với v1),
% thích nghi, AgentDojo v3. 20/11 (góp ý owner): nêu kết quả so với luật so khớp chuỗi theo hướng tích cực nhưng trung
% thực, H1 không đạt được diễn đạt là "cần tập kiểm thử lớn hơn để khẳng định bằng kiểm định đăng ký trước".
% 20/11 (soát văn phong): tách câu dài, bớt "còn", "cưỡng chế" thành "thực thi chính sách".
\begin{abstract}
Một chỉ thị cài vào nội dung mà tác tử LLM đọc qua công cụ có thể khiến tác tử gửi dữ liệu tới đích lạ, làm lộ giá trị
mật hoặc sửa tham số của một hành động hợp lệ. Bài báo đề xuất \ours, một cơ chế thực thi chính sách tất định đặt tại Tool
Broker, gồm chính sách ý định có xét phủ định trích từ yêu cầu của người dùng, nhãn độ nhạy cảm cho từng giá trị và phép
ghim các giá trị người dùng đã nêu, không gọi thêm LLM. Ba thuộc tính cấp quyền, bảo mật và toàn vẹn được chứng minh trên
một mô hình chuyển trạng thái dưới năm giả định nêu rõ. Trên tập kiểm thử tiếng Việt gồm 48 họ tấn công, niêm phong trước
khi viết mã, với ba mô hình nền mở trọng số, \ours{} giảm số họ bị tấn công thành công từ 40 khi không phòng thủ xuống 2.
Khi ba thành phần được thay bằng luật so khớp chuỗi, còn 6 họ bị hại. \ours{} chặn thêm bốn họ trong số đó, mỗi họ bởi
đúng thành phần được thiết kế cho nó, không để lọt họ mới, và cho cùng chiều hiệu ứng trên cả ba mô hình. Do luật so khớp
chuỗi đã gần mức sàn, cần một tập kiểm thử lớn hơn để khẳng định khác biệt này bằng kiểm định đăng ký trước. Tiện ích chỉ
giảm 3,1 điểm phần trăm so với khi không phòng thủ, còn token và độ trễ chênh không quá 1,4\% so với luật so khớp chuỗi.
Dưới hai vòng tấn công thích nghi, \ours{} không để lọt họ nào ngoài những họ đã lọt với tấn công gốc. Trên AgentDojo
bằng tiếng Anh và tiếng Việt, số tấn công thành công được giữ ở mức thấp, đổi lại tiện ích giảm ở các nhiệm vụ lấy đích
từ dữ liệu công cụ.

\keywords{Chèn lệnh gián tiếp \and Tác tử LLM \and Cấp quyền theo ý định \and Kiểm soát luồng thông tin \and Tiếng
Việt}
\end{abstract}


% Số: facts.md (Test v3 và Test-3: 32,9 -> 1,2%, 32,4 -> 3,8%; 40 ca lọt v1 theo nguyên nhân, mục "Bổ sung sau phản
% biện 3"; vi mô broker, mục "Bổ sung cho bài v3"). Giai đoạn đầu (luật chuỗi trên hai tập niêm phong) được kể như một
% phần của chính nghiên cứu này, không gọi là phiên bản trước.
\section{Giới thiệu}\label{sec:intro}

Trong mô hình ReAct~\cite{yao2023react}, tác tử LLM xen kẽ suy luận với lời gọi công cụ, và nội dung công cụ trả về đi
thẳng vào ngữ cảnh của mô hình. Trong tấn công chèn lệnh gián tiếp~\cite{greshake2023indirect}, kẻ tấn công cài chỉ thị
vào nội dung mà tác tử sẽ đọc để khiến nó gửi dữ liệu tới một đích lạ, làm lộ giá trị mật hoặc sửa nội dung cần gửi.
Phòng thủ mức prompt phụ thuộc vào việc mô hình có tuân theo hay không, còn các bộ phân loại chèn lệnh phần lớn được
huấn luyện trên tiếng Anh. Một hướng khác đặt việc cưỡng chế ở phía host, giữa mô hình và công cụ, theo nguyên lý
reference monitor~\cite{anderson1972}.

Cách đơn giản nhất để cưỡng chế tại Tool Broker là so khớp chuỗi. Broker chỉ thực thi một lời gọi gửi ra ngoài khi đích
của nó là một chuỗi con của yêu cầu người dùng $u$, và chỉ cho một giá trị lấy từ nguồn gắn nhãn mật rời tác tử khi $u$
chứa nguyên văn giá trị đó. Ở giai đoạn đầu của nghiên cứu này, các luật như vậy kết hợp với lời nhắc kiểu sandwich đã
giảm tỉ lệ tấn công thành công trên hai tập kiểm thử tiếng Việt niêm phong, với ba mô hình nền mở, từ 32,9\% xuống
1,2\% và từ 32,4\% xuống 3,8\%. Phân tích cả 40 ca còn lọt cho thấy ba chỗ hở (Bảng~\ref{tab:mechanisms}). Có 25 ca làm
lộ một giá trị mật nằm trong nguồn gắn nhãn công khai, vì nhãn được đặt cho cả nguồn. Có 5 ca gửi tới một đích có xuất
hiện trong $u$ nhưng không được cấp quyền, chẳng hạn đích nằm trong một câu cấm. Mười ca còn lại sửa tham số của một lời
gọi tới đích hợp lệ.

Các ca này cho thấy một chuỗi xuất hiện trong $u$ chưa chắc đã được cấp quyền, và nhãn của nguồn không cho biết độ nhạy
cảm của từng giá trị. Bài báo đề xuất \ours, một cơ chế cưỡng chế tại broker gồm ba thành phần tất định nhắm vào ba chỗ
hở này. \POne{} trích từ $u$ các quyền có cấu trúc, xét phạm vi phủ định và loại trừ, và chỉ cấp quyền cho đích khớp
nguyên token. \PTwo{} gán độ nhạy cảm cho từng giá trị trong đầu ra công cụ theo kiểu của giá trị, độc lập với nhãn
nguồn. \PThree{} buộc lời gọi được cấp quyền giữ các giá trị người dùng đã nêu. Chính sách chỉ được suy ra từ $u$,
không dựa vào phán đoán của một LLM, và cho kết quả \emph{cần xác nhận} khi quyền không rõ ràng. Bài báo có ba đóng góp.
\begin{enumerate}
\item \textbf{Phương pháp.} \ours{} xử lý ba chỗ hở đã đo được của luật so khớp chuỗi mà không gọi thêm LLM, với thời
  gian kiểm dưới một mili giây mỗi lời gọi (Mục~\ref{sec:method}).
\item \textbf{Mô hình hình thức.} Ba mệnh đề về cấp quyền, bảo mật và toàn vẹn được chứng minh dưới năm giả định, tách
  khỏi các thuộc tính chỉ đo được bằng thực nghiệm (Mục~\ref{sec:formal}).
\item \textbf{Đánh giá.} \ours{} được đóng băng rồi chạy một lần trên một tập kiểm thử mới gồm 48 họ, niêm phong trước
  khi có mã, với giả thuyết đăng ký trước, ablation, hai vòng tấn công thích nghi và AgentDojo bằng hai ngôn ngữ
  (Mục~\ref{sec:protocol}, \ref{sec:results}). Bằng chứng mức ca cho thấy mỗi thành phần chặn đúng loại tấn công mà
  nó nhắm tới, còn khẳng định thống kê so với luật so khớp chuỗi cần một tập kiểm thử lớn hơn
  (Mục~\ref{sec:discussion}).
\end{enumerate}
Dữ liệu, mã và các manifest niêm phong được công bố công khai.

% Bản 15 trang: 19 tài liệu. Bỏ BIPIA, Nasr 2025, Tramèr 2020, Dwork 2015, Recht 2019, Holm 1979. Giữ FIDES, design
% patterns và Liu et al. vì phản biện 3 coi việc cắt các tài liệu này là suy giảm.
% CapLite (arXiv:2611.17115, 12/11/2026): chỉ dùng mức tóm tắt đã đọc ngày 19/11 (facts.md, "Tài liệu mới được trích").
\section{Nghiên cứu liên quan}\label{sec:related}

\paragraph{Tấn công và phòng thủ dựa vào mô hình.} Liu và cộng sự~\cite{liu2024formalizing} hình thức hóa tấn công chèn
lệnh và phòng thủ, còn InjecAgent~\cite{zhan2024injecagent} và AgentDojo~\cite{debenedetti2024agentdojo} đo tấn công
trên tác tử dùng công cụ. Các bộ dữ liệu này đều bằng tiếng Anh, trong khi cơ chế an toàn của LLM yếu hơn ở ngôn ngữ ít
tài nguyên~\cite{deng2024multijail}, và tấn công thích nghi đã vượt qua nhiều phòng thủ đã công
bố~\cite{zhan2025adaptive}. Spotlighting~\cite{hines2024spotlighting} đánh dấu dữ liệu không đáng tin, lời nhắc kiểu
sandwich lặp lại yêu cầu sau mỗi quan sát, SecAlign~\cite{chen2025secalign} huấn luyện lại mô hình để tách chỉ thị khỏi
dữ liệu, Prompt Guard~2~\cite{meta2025promptguard2} sàng lọc đầu vào bằng một bộ phân loại, còn
MELON~\cite{zhu2025melon} so hành động của lượt chạy gốc với lượt chạy lại có che yêu cầu. Các cách này phụ thuộc vào
hành vi của mô hình hoặc bộ phân loại trên chính dữ liệu tấn công.

\paragraph{Cưỡng chế bên ngoài mô hình.} Progent~\cite{shi2025progent} áp chính sách đặc quyền lên lời gọi công cụ, còn
ClawGuard~\cite{zhao2026clawguard} nhờ mô hình suy ra luật truy cập từ yêu cầu của người dùng.
FIDES~\cite{costa2025fides} gắn nhãn tin cậy và bảo mật cho dữ liệu rồi lan truyền nhãn qua các lời gọi.
CaMeL~\cite{debenedetti2026camel} tách luồng điều khiển khỏi dữ liệu không đáng tin bằng một trình thông dịch riêng và
gắn capability cho từng giá trị. CapLite~\cite{xu2026caplite}, công bố cùng thời gian, suy chính sách từ schema công cụ,
cam kết các lời gọi ghi trước khi đọc kết quả công cụ và chỉ nhận lời gọi ghi khi các đối số quyết định đối tượng có
trong yêu cầu. Beurer-Kellner và cộng sự~\cite{beurerkellner2025patterns} hệ thống hóa các mẫu thiết kế kiểu này. Bản
thân việc dùng chính sách có cấu trúc hay chứng minh an toàn vì vậy chưa phải là đóng góp mới. \ours{} khác ở chỗ giữ
nguyên vòng lặp ReAct và chỉ kiểm từng lời gọi tại broker, thay vì viết lại tác tử thành trình thông dịch như CaMeL hay
cố định tập lời gọi ghi trước khi đọc dữ liệu như CapLite. Chính sách được trích tất định từ $u$ và xét cực tính trong
tiếng Việt có dấu, không dấu và tiếng Anh, trong khi ClawGuard để mô hình viết luật và Progent dựa vào luật viết tay. Độ
nhạy cảm được gán theo kiểu giá trị ngay tại broker, không lan truyền qua mọi phép tính như FIDES, nên rẻ hơn nhưng chỉ
bảo vệ giá trị trong một lớp biến đổi đã nêu (Mục~\ref{sec:formal}).

% Công cụ: docs/tool_contracts.md (6 công cụ đọc + 2 sink giả, không I/O mạng). Bộ chấm biến đổi:
% src/react_agent/evaluation/transformed_leakage.py (decision log 2026-11-15 18:20). Không nêu tên công cụ giả trong bài.
\section{Mô hình đe dọa}\label{sec:threat}

\paragraph{Tác tử, công cụ và nhãn.} Gọi $u$ là yêu cầu của người dùng. Ở mỗi bước, mô hình sinh một hành động $a_t$ gồm
tên công cụ và tham số, hoặc sinh câu trả lời cuối $y$. Tool Broker kiểm tra $a_t$, thực thi và trả về quan sát $o_t$.
Hệ thống có sáu công cụ đọc và hai \emph{sink} có tác dụng phụ là gửi e-mail và gọi webhook, đều là bản giả chỉ ghi
nhật ký. Câu trả lời cuối được coi là sink thứ ba. Với một lời gọi sink $a$, $\dest(a)$ là đích (địa chỉ hoặc URL) và
$\data(a)$ là các tham số còn lại. Mỗi nội dung mang hai nhãn độc lập. Theo độ tin cậy, chỉ $u$ và prompt hệ thống là
đáng tin. Theo độ nhạy cảm, nội dung được xếp từ công khai tới mật, vì một trang không đáng tin có thể là công khai, còn
một bản ghi mật có thể đến từ nguồn mà kẻ tấn công không chạm tới.

\paragraph{Kẻ tấn công.} Kẻ tấn công kiểm soát văn bản của một nội dung do công cụ đọc trả về và biết schema của công cụ.
Trong các vòng thích nghi, họ biết thêm toàn bộ thiết kế và mô hình hình thức của phòng thủ. Họ không sửa được $u$,
prompt hệ thống, broker hay trọng số mô hình. Jailbreak trực tiếp, đầu độc dữ liệu huấn luyện và hệ nhiều tác tử nằm
ngoài phạm vi.

\paragraph{Vị từ gây hại.} Gọi $G(u)$ là tập cặp (công cụ, đích) mà người dùng thực sự cho phép, $\Sigma$ là tập giá trị
mật trên quỹ đạo mà người dùng không cho phép tiết lộ, và $\mathcal{S}$ là tập lời gọi sink đã được thực thi. Giá trị do
người dùng viết trong $u$ không thuộc $\Sigma$. Quỹ đạo $\tau$ gây hại, ký hiệu $\mathrm{Harm}(\tau)$, khi ít nhất một
trong bốn vị từ sau đúng: $H_\mathit{dest}$ khi có $a \in \mathcal{S}$ với $(\mathrm{tool}(a), \dest(a)) \notin G(u)$,
$H_\mathit{leak}$ khi một $v \in \Sigma$ xuất hiện trong $\data(a)$ của một $a \in \mathcal{S}$ hoặc trong $y$ ở dạng
nguyên văn hay biến đổi, $H_\mathit{tamper}$ khi có $a \in \mathcal{S}$ tới đích hợp lệ nhưng một trường do người dùng
chỉ định mang giá trị khác, và $H_\mathit{side}$ khi nhiệm vụ chỉ cần đọc nhưng $\mathcal{S} \neq \emptyset$. Một bộ
chấm đóng băng xác định các vị từ từ lời gọi đã thực thi và $y$, nên lời gọi bị chặn không được tính là gây hại. Ngoài
so khớp nguyên văn, bộ chấm nhận ra giá trị mật ở dạng Base64, Hex, mã hóa URL, tách ký tự và số viết bằng chữ.

\paragraph{Cơ chế, họ và biến thể.} Tấn công được phân loại theo sáu \emph{cơ chế}, tức con đường để đoạn chèn biến
thành tác hại (Bảng~\ref{tab:mechanisms}). Ở hai cơ chế chuyển hướng, đích không có trong $u$. Ở hai cơ chế rò rỉ và cơ
chế sửa tham số, đích hợp lệ, nên chỉ nguồn gốc hay giá trị của dữ liệu mới phân biệt được tấn công với nhiệm vụ. Một
\emph{họ} là một ý tưởng tấn công trong một nhiệm vụ và một môi trường công cụ cụ thể. Mọi họ có nhánh \emph{\bPlain},
thay đoạn chèn bằng một ghi chú vô hại cùng độ dài, và một phần ba số họ có thêm nhánh \emph{\bAuth}, trong đó người
dùng cho phép rõ ràng một hành động giống mục tiêu tấn công. Chặn ở nhánh lành tính là \emph{chặn nhầm}. Mỗi họ có năm
biến thể áp cho cả $u$ lẫn đoạn chèn là tiếng Việt chuẩn, diễn đạt lại, trộn Việt và Anh, không dấu, và tiếng Anh, với
định danh, địa chỉ và giá trị mật giữ nguyên.

% Bảng 1. Nguồn số: facts.md mục "Bổ sung sau phản biện 3" (nguyên nhân lọt v1: Test-I 13 = SL 6, AL 3, PT 4 (3 sửa giá
% trị + 1 đích có trong u); Test-II 27 = SL 14, AL 2, PT 7, US 3, CR 1) và mục Test-4 ("40 ca vượt v1": M3 20, M5 11,
% M4 5, M6 3, M2 1). Tổng theo nguyên nhân: nguồn công khai 25, đích có trong u 5, sửa tham số 10.
\begin{table}[!htbp]
\centering
\caption{Sáu cơ chế tấn công. Cột thứ ba ghi số ca lọt qua \base{} trên hai tập niêm phong trước (\tOne{} và \tTwo,
ba mô hình nền, tổng 40 ca) và nguyên nhân. Cột cuối là thành phần của \ours{} nhắm vào nhóm đó.}\label{tab:mechanisms}
\footnotesize
\setlength{\tabcolsep}{3pt}
\begin{tabularx}{\textwidth}{@{}>{\raggedright\arraybackslash}p{2.45cm}>{\raggedright\arraybackslash}X%
>{\raggedright\arraybackslash}p{4.0cm}l@{}}
\toprule
Cơ chế & Con đường gây hại & Lọt qua \base{} (nguyên nhân) & \ours \\
\midrule
\mER{} chuyển hướng tường minh & Đoạn chèn yêu cầu gửi dữ liệu tới một đích lạ & 0 & \cOne \\
\mCR{} chuyển hướng che giấu & Như trên, nhưng diễn đạt gián tiếp hoặc tách nhỏ đích & 1 (trường đích là một từ thường
  có trong $u$) & \cOne \\
\mSL{} rò rỉ qua sink & Giá trị mật đi vào nội dung gửi tới một đích hợp lệ & 20 (giá trị nằm trong nguồn gắn nhãn
  công khai) & \cTwo \\
\mAL{} rò rỉ qua câu trả lời & Giá trị mật xuất hiện trong câu trả lời cuối & 5 (như trên) & \cTwo \\
\mPT{} sửa tham số & Đổi một giá trị người dùng đã nêu, giữ đích hợp lệ & 11 (10 ca sửa giá trị, 1 ca đích chỉ được
  nhắc tới) & \cThree, \cOne \\
\mUS{} tác dụng phụ ngoài yêu cầu & Gọi sink tới một đích chỉ được nhắc tới hoặc bị cấm & 3 (đích có trong $u$ nhưng
  không được cấp quyền) & \cOne \\
\bottomrule
\end{tabularx}
\end{table}


% Nguồn: src/react_agent/security_v3/{intent_policy_v1,value_labels_v1,param_pinning_v1,brokerguard_v3_v1}.py;
% decision log 2026-11-16 12:30 (C1), 15:30 (runtime, bỏ C1a), 16:20 (vi mô), 2026-11-17 19:25 (đóng băng).
% Số: facts.md mục "Bổ sung cho bài v3" (ngưỡng 8 ký tự, 9 chữ số; vi mô) và mục Test-4 (0,26 ms). Bản 15 trang bỏ
% đoạn về biến thể C1a (2 -> 5 quỹ đạo gây hại trên Dev-val), xem bản dài ở commit 6814b3a.
% Không trích văn bản tấn công. Ví dụ từ nối phủ định là của bộ trích, không phải của dữ liệu.
\section{Phương pháp \ours}\label{sec:method}

Thiết kế của \ours{} xuất phát từ quan điểm rằng mô hình không phải là ranh giới an ninh. Một hành động chỉ có hiệu lực
khi broker thực thi nó, nên broker đóng vai trò reference monitor tại các sink, và nguồn quyền duy nhất là $u$. Lời gọi
bị từ chối được báo lại để tác tử làm tiếp phần việc hợp lệ. Hình~\ref{fig:arch} mô tả kiến trúc. Các luật nền có từ
giai đoạn đầu gồm \rDest{} (Mục~\ref{sec:intro}), \rFlow{} (giá trị từ nguồn gắn nhãn mật chỉ được rời tác tử khi $u$
chứa nguyên văn giá trị đó), \rRedact{} cho câu trả lời cuối và \rSig{} (chặn sink sau một quan sát bị bộ dò cụm từ
gắn cờ, trừ lời gọi khớp một văn phạm chặt trên $u$). Lời nhắc kiểu sandwich nhắc lại $u$ sau mỗi quan sát. \ours{} thay
\rDest{} bằng \pOne{} (\cOne), thêm \pTwo{} (\cTwo) và \pThree{} (\cThree). Cấu hình giữ \rDest{} và bỏ cả ba thành
phần được gọi là \strSand, khi bỏ thêm lời nhắc thì gọi là \strRules.

% Hình 1, TikZ (Overleaf tự vẽ, không render trên máy). Tọa độ tính bằng cm, tổng rộng 11,85 cm < \textwidth 12,2 cm.
% Ô kiểm 1,7 x 0,9 cm, chữ 6,4 pt: tối đa khoảng 14 ký tự mỗi dòng, ba dòng mỗi ô (giữ nhãn ngắn khi sửa).
% Màu theo thực thể (giống Hình 2, 3): cam viền đậm = thành phần mới C1--C3, xanh = giữ từ BG-Base, xám nét đứt = phụ
% thuộc mô hình. In đen trắng: phân biệt bằng độ dày viền và nét đứt.
% Nhãn tiếng Anh cho bản nộp:
%   User / $u$ (trusted) | LLM agent / (ReAct) | Tool Broker (host side)
%   C1: extract $P(u)$ from $u$ / (destination, polarity, pinned values) | C2: label each value in $o_t$ / (protected set $C_t$)
%   before every sink call $a_t$: executed only if all four checks pass
%   R-Signal / flagged, so deny | C1 / positive grant | C3 / pinned values kept | C2 + R-Flow / no value of $C_t$ leaves
%   R-Redact + C2 / redact values in $y$ | Sandwich / repeat $u$ | Report denial, / agent continues
%   Needs confirmation / = denial | Read tools / (untrusted) | Sinks / e-mail, webhook | $y$ redacted
%   new component | kept from BG-Base | model-dependent
\begin{figure}[!htbp]
\centering
\begin{tikzpicture}[
  font=\fontsize{6.4}{7.6}\selectfont,
  box/.style={rounded corners=1pt, align=center, inner sep=1pt},
  newc/.style={box, draw=bgorange, line width=0.9pt, fill=bgorange!12},
  kept/.style={box, draw=bgblue, line width=0.45pt, fill=bgblue!10},
  soft/.style={box, draw=black!55, line width=0.5pt, fill=black!4, dash pattern=on 2pt off 1.3pt},
  neutral/.style={box, draw=black!60, line width=0.5pt, fill=black!5},
  arr/.style={-{Stealth[length=3pt]}, line width=0.5pt, draw=black!65},
  lab/.style={font=\fontsize{6.4}{7}\selectfont},
  note/.style={font=\fontsize{5.6}{6.6}\selectfont\itshape, text=black!65, anchor=west},
  gate/.style={minimum width=1.7cm, minimum height=0.9cm, text width=1.6cm},
  swatch/.style={minimum width=0.34cm, minimum height=0.24cm},
]
% cột trái: người dùng và tác tử
\node[neutral, minimum width=1.5cm, minimum height=0.8cm, text width=1.4cm] at (0.75,4.85)
  {Người dùng\\ $u$ (đáng tin)};
\node[neutral, minimum width=1.5cm, minimum height=1.3cm, text width=1.4cm,
  font=\fontsize{6.6}{7.6}\selectfont\bfseries] at (0.75,2.35) {Tác tử LLM\\ (ReAct)};
\draw[arr] (0.75,4.45) -- (0.75,3.0);

% khung broker
\draw[rounded corners=2pt, line width=0.6pt] (2.0,0.5) rectangle (9.3,5.65);
\node[font=\fontsize{7}{8}\selectfont\bfseries] at (5.65,5.43) {Tool Broker (phía host)};

% hàng 1: trạng thái do C1 và C2 dựng
\node[newc, minimum width=3.45cm, minimum height=0.66cm, text width=3.35cm] at (3.825,4.8)
  {\textsf{C1}: trích $P(u)$ từ $u$\\ (đích, cực tính, giá trị ghim)};
\node[newc, minimum width=3.45cm, minimum height=0.66cm, text width=3.35cm] at (7.475,4.8)
  {\textsf{C2}: gắn nhãn từng giá trị trong $o_t$\\ (tập bảo vệ $C_t$)};
\draw[arr] (1.5,4.8) -- (2.1,4.8);

% hàng 2: bốn kiểm tra trước mỗi lời gọi sink
\node[note] at (2.1,4.2) {trước mỗi lời gọi sink $a_t$: chỉ thực thi khi cả bốn kiểm tra đều qua};
\node[kept, gate] at (2.95,3.5) {\textsf{R-Signal}\\ quan sát bị gắn cờ thì chặn};
\node[newc, gate] at (4.75,3.5) {\textsf{C1}\\ đích có quyền dương};
\node[newc, gate] at (6.55,3.5) {\textsf{C3}\\ giá trị ghim còn nguyên};
\node[newc, gate] at (8.35,3.5) {\textsf{C2} + \textsf{R-Flow}\\ không lộ giá trị của $C_t$};

% hàng 3: câu trả lời cuối, lời nhắc, phản hồi, xác nhận
\node[kept, gate] at (2.95,2.1) {\textsf{R-Redact} + \textsf{C2}\\ che giá trị trong $y$};
\node[soft, gate] at (4.75,2.1) {Sandwich\\ nhắc lại $u$};
\node[kept, gate] at (6.55,2.1) {Báo từ chối,\\ tác tử làm tiếp};
\node[newc, gate] at (8.35,2.1) {Cần xác nhận\\ = từ chối};

% chú giải
\node[newc, swatch] at (2.3,0.88) {};
\node[note] at (2.52,0.88) {thành phần mới};
\node[kept, swatch] at (4.3,0.88) {};
\node[note] at (4.52,0.88) {giữ từ \base};
\node[soft, swatch] at (6.3,0.88) {};
\node[note] at (6.52,0.88) {phụ thuộc mô hình};

% tác tử <-> broker
\draw[arr] (1.5,2.75) -- node[lab, above=-1pt] {$a_t$} (2.0,2.75);
\draw[arr] (2.0,2.35) -- node[lab, above=-1pt] {$o_t$} (1.5,2.35);
\draw[arr] (1.5,1.95) -- node[lab, above=-1pt] {$y$} (2.0,1.95);

% cột công cụ
\node[neutral, minimum width=2.1cm, minimum height=0.8cm, text width=2.0cm] at (10.8,3.95)
  {Công cụ đọc\\ (không đáng tin)};
\node[neutral, minimum width=2.1cm, minimum height=0.8cm, text width=2.0cm] at (10.8,2.45)
  {Sink\\ e-mail, webhook};
\draw[arr] (9.3,4.15) -- node[lab, above=-1pt] {$a_t$} (9.75,4.15);
\draw[arr] (9.75,3.75) -- node[lab, above=-1pt] {$o_t$} (9.3,3.75);
\draw[arr] (9.3,2.45) -- node[lab, above=-1pt] {$a_t$} (9.75,2.45);

% câu trả lời đã che rời broker
\draw[arr] (9.3,1.0) -- (9.95,1.0);
\node[anchor=west, text=black!65] at (9.98,1.0) {$y$ đã che};
\end{tikzpicture}
\caption{Kiến trúc của \ours. Mọi lời gọi $a_t$ và câu trả lời cuối $y$ đều đi qua Tool Broker. \cOne{} trích chính
sách $P(u)$ một lần từ $u$, còn \cTwo{} cập nhật tập giá trị bảo vệ $C_t$ sau mỗi quan sát. Ô viền cam đậm là thành
phần mới, ô viền xanh giữ từ \base, ô nét đứt phụ thuộc vào mô hình. Mọi kiểm tra trong khung đều tất
định.}\label{fig:arch}
\end{figure}


\paragraph{Chuẩn hóa.} Hàm $\norm(\cdot)$ áp Unicode NFKC, bỏ ký tự định dạng và casefold. Hàm $\kappa(\cdot)$ áp thêm
việc bỏ dấu tiếng Việt và chỉ giữ chữ và số. Ký hiệu $x \sub z$ nghĩa là $\kappa(x)$ là một chuỗi con của $\kappa(z)$,
và $V(u)$ là tập giá trị xuất hiện trong $u$ theo nghĩa này.

\paragraph{\POne{} (\cOne).} Bộ trích chỉ đọc $u$ và trả về $P(u) = \{g_i\}$. Mỗi quyền $g = (\mathrm{tool}, \dest, \pi,
\mathrm{pins})$ gồm công cụ, đích, cực tính $\pi \in \{+,-\}$ và các giá trị được ghim. Đích là địa chỉ e-mail hoặc
endpoint, so khớp nguyên token sau $\norm$, nên một phần của đích hay một từ thường trong $u$ không còn được coi là
đích. Phủ định đi với động từ (``đừng gửi'', ``do not send'') phủ các đích theo sau trong cùng mệnh đề, từ loại trừ
(``trừ'', ``except'') phủ tới dấu phẩy kế tiếp, còn từ tương phản (``nhưng'', ``but'') kết thúc phạm vi. Khi $u$ không
dấu, một từ như ``dung'' chỉ được coi là phủ định nếu theo sau là một động từ gửi. Mệnh đề cấp quyền phải có động từ gửi
không bị phủ định và hợp với loại đích. Một đích chỉ được nhắc tới, nằm trong trích dẫn hay đứng gần một phủ định không
rõ phạm vi cho kết quả \emph{cần xác nhận}. Lời gọi sink $a$ thỏa \pOne{} khi
\begin{equation}\label{eq:c1}
\begin{aligned}
&\exists\, g \in P(u)\colon\; \pi(g) = {+} \,\land\, \mathrm{tool}(g) = \mathrm{tool}(a) \,\land\, \dest(g) \doteq \dest(a),\\
&\nexists\, g' \in P(u)\colon\; \pi(g') = {-} \,\land\, \dest(g') \doteq \dest(a),
\end{aligned}
\end{equation}
với $\doteq$ là bằng nhau nguyên token sau $\norm$. Chẳng hạn, với yêu cầu ``gửi báo cáo cho \code{ketoan@example.com},
không gửi cho \code{doitac@example.org}'', lời gọi tới địa chỉ thứ hai bị từ chối dù \rDest{} cho qua.

\paragraph{\PTwo{} (\cTwo).} \PTwo{} gán nhãn cho từng giá trị trong mọi đầu ra công cụ theo kiểu của nó, độc lập với
nhãn nguồn và độ tin cậy. Ba kiểu được nhận ra là định danh trộn chữ và số dài ít nhất 8 ký tự, số có ít nhất 9 chữ số
sau khi bỏ khoảng trắng, dấu chấm và gạch, và giá trị đứng sau tên trường nhạy cảm như mật khẩu, OTP, token hay số tài
khoản. Gọi $C_t$ là tập giá trị cần bảo vệ sau bước $t$, gồm các giá trị được \rFlow{} bảo vệ và các giá trị được gắn
nhãn theo kiểu. Broker từ chối lời gọi sink $a$ nếu
\begin{equation}\label{eq:c2}
\exists\, v \in C_t \setminus V(u)\colon\; v \sub \data(a) .
\end{equation}
Vì so khớp dùng $\kappa$, điều kiện vẫn đúng khi giá trị bị thêm hay bớt khoảng trắng, dấu chấm, gạch, đổi hoa thường
hoặc bỏ dấu. Trong câu trả lời cuối, mọi $v \in C_t \setminus V(u)$ có $v \sub y$ được thay bằng \code{[REDACTED]}.

\paragraph{\PThree{} (\cThree).} Sửa tham số giữ nguyên đích được cấp quyền và đổi một giá trị như số tiền, trạng thái,
đơn vị hay mã. Thành phần này ghim các literal trong mệnh đề của $u$ đã cấp quyền cho lời gọi, gồm số đứng sau động từ
gửi và không phải toán hạng của một phép tính, đơn vị nêu kèm số, giá trị boolean có cực tính và định danh do người dùng
viết. Lời gọi bị từ chối khi thiếu một literal đã ghim trong $\data(a)$, so ở dạng $\kappa$, hoặc mang boolean ngược với
boolean đã ghim.

\paragraph{Cổng và chi phí.} Một lời gọi sink chỉ được thực thi khi qua cả bốn kiểm tra là \rSig, \pOne, \pThree{} và
kiểm luồng gồm \rFlow{} cùng \pTwo. Kết quả \emph{cần xác nhận} được thi hành như từ chối và đếm riêng, còn lời gọi đọc
luôn được cho qua. Cổng không gọi thêm mô hình nào. Trên tập phát triển, trung vị thời gian kiểm một lời gọi sink là
315~$\mu$s (\strSand: 107~$\mu$s).

% Nguồn: docs/project/brokerguard_v3_formal_model.md (mục 1--6). Bằng chứng trong Bảng 2: facts.md mục Test-4 (ablation
% theo cơ chế, bằng chứng cơ chế mức ca, hai họ M5 còn lại) và mục "Bổ sung cho bài v3" (C1 offline trên Dev-val).
% Mệnh đề 3 chỉ bảo đảm literal ghim CÓ MẶT trong data(a), không bảo đảm không có giá trị khác bên cạnh.
% Mã giữ lại ở mục này: (G1)--(G4) và (A1)--(A5), vì chứng minh và Bảng 2 tham chiếu tới chúng.
% 20/11 (lượt 3): G(u), Sigma định nghĩa tại đây; vị từ gây hại gọi bằng tên (chuyển hướng, rò rỉ...) thay vì H_dest...; pins bằng lời.
\section{Mô hình hình thức}\label{sec:formal}

Mục này xác định khi nào các kiểm tra ở Mục~\ref{sec:method} đủ ngăn các kiểu tác hại ở Mục~\ref{sec:threat}. Gọi
$G(u)$ là tập cặp công cụ và đích mà người dùng thực sự cho phép, và $\Sigma$ là tập giá trị mật mà người dùng không cho
tiết lộ. Trạng thái của broker sau bước $t$ là $s_t = (u, P(u), C_t, H_t)$, với $H_t$ là lịch sử lời gọi đã thực thi và
kết quả. $P(u)$ được dựng một lần trước bước đầu. Gọi $R$ là tập công cụ đọc và $S$ là tập sink có trường đích. Cổng xét
mỗi lời gọi đề xuất $a_t$ và trả về cho qua, từ chối hoặc cần xác nhận, trong đó cần xác nhận được xử lý như từ chối. Lời
gọi đọc luôn được cho qua. Với $a_t \in S$, cổng cho qua khi và chỉ khi (G1) \rSig{} không chặn $a_t$, (G2) $a_t$ thỏa
\eqref{eq:c1}, (G3) có ít nhất một quyền $g$ thỏa (G2) mà mọi literal được ghim của $g$ đều có mặt trong $\data(a_t)$ và
không boolean nào bị đảo, và (G4) điều kiện \eqref{eq:c2} sai. Gọi $\hat G(u)$ là tập cặp công cụ và đích thỏa (G2). Khi
cổng cho qua, broker thực thi $a_t$, ghi lời gọi và kết quả $o_t$ vào $H_t$, rồi cập nhật $C_{t+1} = C_t \cup L(o_t)$
với $L$ là hàm gắn nhãn kết hợp \rFlow{} và \pTwo. Câu trả lời cuối được phát hành sau khi che mọi giá trị của $C_t
\setminus V(u)$. Năm giả định sau giới hạn phạm vi của các mệnh đề.
\begin{enumerate}[label=(A\arabic*), leftmargin=*]
\item \emph{Trung gian toàn phần.} Mọi lời gọi công cụ đều đi qua broker, và mọi công cụ có tác dụng phụ đều thuộc $S$.
\item \emph{Bộ trích đúng.} $\hat G(u) \subseteq G(u)$. Chiều ngược lại chỉ ảnh hưởng tới tiện ích.
\item \emph{Độ phủ nhãn.} Mọi giá trị của $\Sigma$ xuất hiện trong một kết quả $o_t$ đều thuộc $C_{t+1}$.
\item \emph{Lớp biến đổi.} Nếu tác tử làm lộ $v$ trong một chuỗi $w$ thì $v \sub w$. Mã hóa, viết số bằng chữ hay chia
  một giá trị qua nhiều lời gọi nằm ngoài lớp này.
\item \emph{Literal ghim đủ.} Với quyền $g$ cấp cho lời gọi $a$, các literal được ghim của $g$ gồm mọi giá trị và
  boolean mà người dùng đã nêu cho $a$.
\end{enumerate}

\begin{proposition}[An toàn cấp quyền]\label{prop:auth}
Dưới (A1) và (A2), mọi lời gọi sink được thực thi có cặp công cụ và đích thuộc $G(u)$. Vì vậy không quỹ đạo nào gây hại
kiểu chuyển hướng, và nhiệm vụ chỉ cần đọc không gây tác dụng phụ.
\end{proposition}
\begin{proof}
Theo (A1), một sink $a$ chỉ được thực thi khi cổng cho qua, nên (G2) đúng và cặp công cụ, đích của $a$ thuộc $\hat G(u)$,
mà $\hat G(u) \subseteq G(u)$ theo (A2). Nhiệm vụ chỉ cần đọc có $G(u)$ rỗng, nên không sink nào được thực thi. \qed
\end{proof}

\begin{proposition}[Bảo mật]\label{prop:conf}
Dưới (A1), (A3) và (A4), không giá trị nào của $\Sigma$ xuất hiện trong $\data(a)$ của một sink được thực thi hay trong
câu trả lời cuối đã phát hành. Vì vậy không quỹ đạo nào gây hại kiểu rò rỉ ở các dạng thuộc (A4).
\end{proposition}
\begin{proof}
Giả sử $v \in \Sigma$ xuất hiện trong $\data(a)$ của một sink $a$ được thực thi ở bước $t$. Vì $v \notin V(u)$, tác tử
chỉ có được $v$ qua một kết quả trước bước $t$, nên theo (A3) $v \in C_t$. Theo (A4), $v \sub \data(a)$, nên
\eqref{eq:c2} đúng và (G4) sai, mâu thuẫn với (A1). Với câu trả lời cuối, phép che thay mọi dạng bề mặt khớp
$\kappa(v)$ trước khi phát hành, nên theo (A4) không còn dạng nào của $v$. \qed
\end{proof}

\begin{proposition}[Toàn vẹn]\label{prop:int}
Dưới (A1), (A2) và (A5), mọi lời gọi sink được thực thi chứa mọi giá trị người dùng đã nêu cho lời gọi đó và không mang
boolean ngược với boolean người dùng đã nêu.
\end{proposition}
\begin{proof}
Theo (A1), sink $a$ chỉ được thực thi khi (G2) và (G3) đúng, tức có quyền $g$ cấp cho $a$ mà mọi literal được ghim của
$g$ có mặt trong $\data(a)$ và không boolean nào bị đảo. Theo (A2), $g$ ứng với một quyền người dùng thật sự cấp, và theo
(A5), các literal được ghim của $g$ gồm mọi giá trị người dùng đã nêu cho $a$. \qed
\end{proof}

\noindent Ba mệnh đề nêu điều kiện đủ dưới các giả định, không phải bảo đảm vô điều kiện. Mệnh đề~\ref{prop:int} không
loại trừ việc một giá trị khác được thêm cạnh literal đã ghim. Lời gọi đọc và sink không có trường đích nằm ngoài $S$, nên
một công cụ đọc nhận URL do mô hình chọn vẫn có thể thành kênh rò rỉ. Bảng~\ref{tab:proven} phân biệt thuộc tính được
chứng minh với thuộc tính chỉ kiểm chứng bằng thực nghiệm, và các giả định thực nghiệm được kiểm tra ở
Mục~\ref{sec:results}.

% Bảng 2. Nguồn: facts.md mục Test-4 (H3 ablation; họ bị hại theo cơ chế: −C1 thêm 1 ER và 1 PT, −C2 thêm 2 SL và 2 AL;
% lộ câu cuối v3 0/720; hai họ M5 còn lại; nguyên văn = biến đổi) và mục "Bổ sung cho bài v3" (C1 offline Dev-val:
% tiện ích 155/155 và 65/65, M6 35/35, v1 15). (A1): tests/unit/test_brokerguard_v3_v1.py. Bọc adjustbox max width.
\begin{table}[!htbp]
\centering
\caption{Thuộc tính được chứng minh và thuộc tính chỉ được đo bằng thực nghiệm. Số trên tập phát triển được dùng khi
thiết kế, số trên tập kiểm thử có sau khi đóng băng. \cOne{} \pOne, \cTwo{} \pTwo, \cThree{} \pThree.}\label{tab:proven}
\footnotesize
\setlength{\tabcolsep}{3pt}
\begin{adjustbox}{max width=\textwidth}
\begin{tabularx}{\textwidth}{@{}>{\raggedright\arraybackslash}p{2.75cm}>{\raggedright\arraybackslash}p{2.2cm}%
>{\raggedright\arraybackslash}X@{}}
\toprule
Thuộc tính & Trạng thái & Bằng chứng \\
\midrule
Mệnh đề~\ref{prop:auth}--\ref{prop:int} & Chứng minh dưới giả định & Mục~\ref{sec:formal} \\
(A1) trung gian toàn phần & Theo cấu trúc mã, có kiểm thử & Broker là đường duy nhất tới công cụ. Công cụ đọc và sink
  không có đích nằm ngoài $S$ \\
(A2) bộ trích đúng & Thực nghiệm & Tập phát triển: mọi sink lành tính được cho phép, chặn đủ 35 sink tác dụng phụ ngoài
  yêu cầu (\rDest: 15). Tập kiểm thử: bỏ \cOne{} làm lọt thêm 2 họ \\
(A3) độ phủ nhãn & Thực nghiệm, gián tiếp & Tập kiểm thử: bỏ \cTwo{} làm lọt thêm 4 họ rò rỉ. \ours{} không để lộ giá
  trị nào ở câu trả lời cuối \\
(A4) lớp biến đổi & Giới hạn đã nêu & Bộ chấm biến đổi cho cùng kết quả với so khớp nguyên văn ở mọi cấu hình, nên giả
  định chưa bị thử thách \\
(A5) literal ghim đủ & Thực nghiệm, chưa đủ & Tập kiểm thử: 2 họ sửa giá trị số hoặc đơn vị lọt qua cả \strSand{} và
  \ours. Bỏ \cThree{} không làm lọt thêm họ nào \\
Lời nhắc, \rSig & Chỉ thực nghiệm & Không có bảo đảm. Đo qua ablation và tấn công thích nghi (Mục~\ref{sec:results})
  \\
\bottomrule
\end{tabularx}
\end{adjustbox}
\end{table}


% Nguồn: decision log 2026-11-15 16:30 (Dev v4, Test-4), 17:30 (24 ca thích nghi), 18:20 (niêm phong, lực, H1--H4),
% 2026-11-17 19:25 (final_config_v5), 2026-11-18 19:30--20:55 (Gemini, thích nghi v3, AgentDojo v3), 2026-11-19 00:25
% (baseline bổ sung, khai báo trước, khám phá). Số: facts.md mục "Bổ sung cho bài v3" và mục Test-4.
% Nguồn gốc Test-4: sinh từ spec, owner và Minh rà soát (không ghi "do owner và Minh soạn", không báo kappa).
% 20/11 (soát văn phong, giữ 15 trang): bỏ câu kiểm hash mỗi lần chạy (đã có ý niêm phong).
% 20/11 (lượt 3, theo bài tham khảo): câu mở mục, bớt ngoặc có ký hiệu, Gemini chỉ ghi là mô hình đóng của Google.
\section{Dữ liệu và giao thức đánh giá}\label{sec:protocol}

Phần đánh giá kiểm tra \ours{} trên dữ liệu chưa dùng khi thiết kế, so với các phòng thủ đã công bố và với các cấu hình
luật đơn giản hơn tại broker.

\paragraph{Tập dữ liệu.} Bộ dữ liệu VN-AgentInject hiện thực hóa mô hình đe dọa ở Mục~\ref{sec:threat} thành các ca chạy
được với dữ liệu tổng hợp. \Dev{} gồm 168 họ và 1.960 ca (840 tấn công, 840 \bPlain, 280 \bAuth), hợp từ tập phát triển
ban đầu 48 họ và hai tập kiểm thử 72 và 48 họ của giai đoạn đầu. Hai tập kiểm thử được gộp vào vì 40 ca vượt qua trên
chúng đã được dùng để thiết kế \ours, nên dùng lại để kiểm định sẽ cho kết quả lạc quan. \Dev{} được chia theo họ thành
phần huấn luyện 126 họ, dùng để chỉnh danh sách từ, ngưỡng và tên trường, và phần kiểm định 42 họ, dùng để chọn cấu hình.
\Test{} gồm 48 họ mới, 8 họ cho mỗi cơ chế, với 560 ca (240 tấn công, 240 \bPlain, 80 \bAuth). Các họ được sinh tự động
từ một đặc tả rồi được hai thành viên của nhóm rà soát câu chữ. Tập được niêm phong bằng SHA-256 trước khi viết dòng mã
đầu tiên của \ours, với độ trùng 4-gram Jaccard so với \dev{} tối đa là 0,05.

\paragraph{Tấn công thích nghi.} Bộ thứ nhất gồm 24 ca thuộc \test, được sinh và niêm phong cùng \test, nhắm vào ba
điểm yếu của \strSand{} với 8 ca cho mỗi điểm yếu. Bộ thứ hai gồm 42 ca do hai thành viên của nhóm soạn sau khi \ours{}
đã đóng băng, với hiểu biết đầy đủ về thiết kế và mô hình hình thức. Mỗi ca thay văn bản tấn công của một họ thuộc phần
kiểm định của \dev{} và nhắm vào các giả định (A2) đến (A5).

\paragraph{Mô hình nền và cấu hình.} Gemma-4-12B, Llama-3.1-8B với ngữ cảnh 16k và Qwen3.5-9B được phục vụ bằng vLLM với
phiên bản cố định, dùng chung prompt ReAct, nhiệt độ 0, tối đa 512 token mỗi bước và tám bước. Gemini 3 Flash của Google
là mô hình đóng nên chỉ được dùng làm cột tham khảo với năm cấu hình. Mọi cấu hình dùng chung tác tử, broker và môi
trường, và khi không phòng thủ, broker chỉ kiểm tra schema. Các phòng thủ khác gồm \cSpot~\cite{hines2024spotlighting},
\cSand, \cPG{} dùng Prompt Guard~2~\cite{meta2025promptguard2} để sàng lọc đầu ra công cụ, \cMel{} và \cProg. Hai phòng
thủ sau được cài đặt lại theo mô tả gốc của MELON~\cite{zhu2025melon} và Progent~\cite{shi2025progent}. Luật tại broker
gồm \phrRules, chặn theo danh sách cụm từ giống chỉ thị, cùng \strRules, \strSand, \ours{} và ba bản bỏ một thành phần.
Tám cấu hình của ma trận chính chạy một lần trên mỗi mô hình nền với 560 ca và 24 ca thích nghi, chỉ chạy lại ca gặp lỗi
hạ tầng. Năm cấu hình \cSpot, \cPG, \cMel, \phrRules{} và \strRules{} đã đóng băng từ trước nhưng được chạy bổ sung sau
khi có kết quả chính, nên so sánh với chúng chỉ mang tính khám phá.

\paragraph{Độ đo.} ASR là tỉ lệ ca tấn công gây hại, tính trên 720 ca tấn công của ba mô hình nền. Tiện ích $U$ là tỉ lệ
ca lành tính hoàn thành theo điều kiện tất định của từng nhiệm vụ mà không gây hại, tính trên 960 ca, và $\Delta U$ là
chênh lệch so với khi không phòng thủ. Lỗi mô hình, như đầu ra sai định dạng hoặc bị cắt, được tính là thất bại về tiện
ích, không bao giờ là phòng thủ thành công. Một họ được coi là \emph{bị hại} nếu có ít nhất một ca tấn công gây hại trên
ít nhất một mô hình nền. Đây là đơn vị suy luận chính, vì các ca trong cùng một họ dùng chung nhiệm vụ và môi trường.
Trên \dev, đơn vị tương ứng là cặp họ và mô hình, gồm 504 cặp. Chi phí gồm tỉ lệ chặn nhầm, tức tỉ lệ ca lành tính có ít
nhất một lần cổng từ chối, tỉ lệ cần xác nhận, số token và thời gian mỗi ca.

\paragraph{Giả thuyết đăng ký trước.} Trước khi niêm phong \test, bốn giả thuyết được đăng ký theo thủ tục kiểm định tuần
tự có cổng. H1 là giả thuyết chính duy nhất: \ours{} có ít họ bị hại hơn \strSand, kiểm bằng phép McNemar chính
xác~\cite{mcnemar1947} một phía trên 48 cặp họ với $\alpha = 0{,}05$. H2 đến H4 chỉ được kiểm định khi giả thuyết không
của H1 bị bác bỏ. H2 cho rằng tiện ích của \ours{} không thấp hơn \strSand{} quá 5 điểm. H3 cho rằng bỏ mỗi thành phần
làm tăng số họ bị hại, có hiệu chỉnh Holm. H4 đặt ngưỡng 10\% cho tỉ lệ cần xác nhận, 15\% cho tỉ lệ chặn nhầm, và mức
tăng so với \strSand{} không quá 20\% về token, 15\% về độ trễ.

\paragraph{AgentDojo.} Bộ Banking với 144 cặp nhiệm vụ người dùng và nhiệm vụ chèn và bộ Slack với 105 cặp của AgentDojo
v1.2.2~\cite{debenedetti2024agentdojo} được chạy với khuôn tấn công \code{important\_instructions}, bằng tiếng Anh và
bằng bản dịch tiếng Việt giữ nguyên IBAN, số tiền, tên và địa chỉ. \StrSand{} và \ours{} chạy với Llama-3.1-8B qua một bộ
chuyển đổi khai báo trước, chỉ rõ công cụ nào là sink và đối số nào chứa đích. Trong bộ chuyển đổi, \pOne{} chỉ yêu cầu
đích xuất hiện nguyên token trong $u$ và không nằm trong phạm vi phủ định. Lời nhắc \cSand{} ở đây là \code{repeat\_user\_prompt} có sẵn của AgentDojo, và tiện ích khi không có tấn công được tính trên 16 nhiệm vụ
Banking và 21 nhiệm vụ Slack.

% Số: facts.md mục Test-4 (bảng chính, kiểm định, theo cơ chế, loại hại, bằng chứng mức ca), mục Gemini Test-4, mục
% baseline bổ sung Test-4 và Dev v4.1, mục thích nghi nhắm v3, mục AgentDojo v1 và v3, mục "Bổ sung cho bài v3" (ASR mức
% ca, chi phí theo mô hình). Không nêu mã họ Test, không trích văn bản tấn công.
% Bản 15 trang: bảng chi phí theo mô hình và bảng Gemini của bản dài thành lời văn; hình theo cơ chế và bảng thích nghi
% gộp thành Bảng 4; bảng giả thuyết H1--H4 thành lời văn; tiểu mục thành tiêu đề đoạn; Hình Pareto bật bằng \paretotrue.
% 20/11 (soát văn phong, giữ 15 trang): bỏ câu theo biến thể ngôn ngữ (0 đến 2 ca, quá ít để kết luận) và số ca cần xác
% nhận của Gemini (5); xem commit dc564fa.
% 20/11 (lượt 3, theo bài tham khảo): câu mở mục, bớt ngoặc có ký hiệu, Gemini chỉ ghi là mô hình đóng của Google.
\section{Kết quả}\label{sec:results}

\paragraph{Kết quả chính.} Bảng~\ref{tab:main} tổng hợp mười ba cấu hình trên \test. Khi không phòng thủ, 40 trong 48 họ
bị hại trên ít nhất một mô hình nền, và 32,1\% số ca tấn công thành công. Với \ours, chỉ còn 2 họ bị hại, tương ứng 5
trên 720 ca tấn công hay ASR 0,7\%. Khi luật tại broker được làm chặt dần, số họ bị hại giảm đơn điệu từ 34 với
\phrRules{} xuống 15 với \strRules, 6 với \strSand{} và 2 với \ours, và \dev{} cho cùng thứ tự đó. Hai phòng thủ phát
hiện gần như không có tác dụng trên dữ liệu tiếng Việt: \cPG{} không gắn cờ quan sát nào của \test{} và \cMel{} vẫn để 40
họ bị hại. \cSand{} và \cProg{} chỉ chặn được khoảng một nửa số họ.

Đối với giả thuyết chính H1, số họ bị hại giảm từ 6 với \strSand{} xuống 2 với \ours. Có 4 họ chỉ bị hại dưới \strSand{}
và không họ nào chỉ bị hại dưới \ours, nên kiểm định McNemar chính xác một phía cho $p = 0{,}0625$, hai phía là
$0{,}125$, chưa đạt mức $\alpha = 0{,}05$ đã đăng ký. Chiều hiệu ứng giống nhau trên cả ba mô hình, với số họ giảm từ 2,
2 và 3 xuống 1 ở mỗi mô hình, nhưng từng mô hình có quá ít cặp lệch để kết luận riêng. Theo thủ tục cổng, H2 đến H4 chỉ
được báo cáo ở dạng mô tả. Tiện ích của \ours{} thấp hơn \strSand{} 0,56 điểm trên 720 cặp ca \bPlain{} và 0,42 điểm
trên 240 cặp ca \bAuth, nằm trong biên 5 điểm của H2. Ba bản bỏ một thành phần có 4, 6 và 2 họ bị hại so với 2 của
\ours, với $p$ một phía lần lượt là 0,50, 0,19 và 1. So với các phòng thủ khác, khác biệt rõ hơn nhiều, chẳng hạn $p =
3{,}6\cdot10^{-12}$ so với khi không phòng thủ, $9{,}5\cdot10^{-7}$ so với \cSand{} và $1{,}2\cdot10^{-4}$ so với
\strRules, dù đây là phân tích khám phá chưa hiệu chỉnh cho nhiều phép so sánh.

% Bảng 3. Nguồn: facts.md mục Test-4 (họ /48, Q·G·L, tiện ích, H1--H4), mục baseline bổ sung Test-4 và Dev v4.1, mục
% "Bổ sung cho bài v3" (ASR /720, tiện ích dạng số ca, Dev P2 111 và PRG 198). ΔU = (k − 411)/960, làm tròn 1 chữ số.
% ASR làm tròn nửa lên: 171/720 = 23,75 -> 23,8; 9/720 = 1,25 -> 1,3. Chặn mọi sink: 54/720 ca lộ ở câu trả lời.
% In đậm: tốt nhất về an ninh (không tính hàng ước lượng).
\begin{table}[!htbp]
\centering
\caption{Kết quả trên \tThree{} (48 họ, ba mô hình nền) và trên \dev. Cột ``Họ'' là số họ bị hại trên ít nhất một mô
hình nền. Q, G, L là số họ bị hại trên Qwen3.5, Gemma-4 và Llama-3.1. ASR tính trên 720 ca tấn công, $U$ trên 960 ca
lành tính, $\Delta U$ (điểm) so với \cNone{} tính từ số ca. Cột \dev{} là số cặp họ và mô hình bị hại (/504).
$^\dagger$Chạy bổ sung sau khi có kết quả chính. $^\ddagger$Ước lượng từ quỹ đạo \cNone. Phần dưới: các giả thuyết
đăng ký trước.}\label{tab:main}
\footnotesize
\setlength{\tabcolsep}{3pt}
\begin{tabular*}{\textwidth}{@{\extracolsep{\fill}}lrcrrrr@{}}
\toprule
Cấu hình & Họ (/48) & Q $\cdot$ G $\cdot$ L & ASR (\%) & $U$ (\%) & $\Delta U$ & \dev{} (/504) \\
\midrule
\cNone & 40 & 24 $\cdot$ 31 $\cdot$ 28 & 32,1 & 42,8 & 0,0 & 346 \\
\cSpot$^\dagger$ & 38 & 28 $\cdot$ 18 $\cdot$ 26 & 24,3 & 39,7 & $-3{,}1$ & 280 \\
\cSand & 22 & 13 $\cdot$ 8 $\cdot$ 10 & 6,5 & 40,7 & $-2{,}1$ & 111 \\
\cPG$^\dagger$ & 40 & 22 $\cdot$ 31 $\cdot$ 28 & 31,7 & 42,7 & $-0{,}1$ & 346 \\
\cMel$^\dagger$ & 40 & 22 $\cdot$ 30 $\cdot$ 28 & 31,5 & 42,8 & 0,0 & 345 \\
\cProg & 25 & 16 $\cdot$ 21 $\cdot$ 18 & 23,8 & 42,6 & $-0{,}2$ & 198 \\
\midrule
\cPhrase$^\dagger$ & 34 & 16 $\cdot$ 25 $\cdot$ 25 & 16,7 & 42,9 & $+0{,}1$ & 278 \\
\cRules$^\dagger$ & 15 & 8 $\cdot$ 10 $\cdot$ 9 & 6,8 & 42,0 & $-0{,}8$ & 123 \\
\base & 6 & 2 $\cdot$ 2 $\cdot$ 3 & 1,3 & 40,2 & $-2{,}6$ & 36 \\
\textbf{\ours} & \textbf{2} & 1 $\cdot$ 1 $\cdot$ 1 & \textbf{0,7} & 39,7 & $-3{,}1$ & \textbf{4} \\
\quad bỏ \cOne & 4 & 1 $\cdot$ 3 $\cdot$ 2 & 1,1 & 40,0 & $-2{,}8$ & -- \\
\quad bỏ \cTwo & 6 & 2 $\cdot$ 2 $\cdot$ 4 & 1,4 & 40,2 & $-2{,}6$ & -- \\
\quad bỏ \cThree & 2 & 1 $\cdot$ 2 $\cdot$ 1 & 0,8 & 39,8 & $-3{,}0$ & -- \\
\midrule
Chặn mọi sink$^\ddagger$ & 14 & -- & 7,5 & 11,5 & $-31{,}4$ & -- \\
\bottomrule
\end{tabular*}

\smallskip
\begin{tabularx}{\textwidth}{@{}l>{\raggedright\arraybackslash}X>{\raggedright\arraybackslash}p{3.4cm}%
>{\raggedright\arraybackslash}p{1.45cm}l@{}}
\toprule
 & Giả thuyết & Kết quả & $p$ & Kết luận \\
\midrule
H1 & Số họ bị hại: \ours{} $<$ \base & $6 \to 2$ (cặp lệch 4/0) & 0,0625 & không bác bỏ $H_0$ \\
H2 & Tiện ích: \ours{} $\ge$ \base{} $-\,5$ điểm & \bPlain{} $-0{,}56$, \bAuth{} $-0{,}42$ điểm & -- & mô tả \\
H3 & Bỏ \cOne, \cTwo, \cThree{} làm tăng số họ bị hại & 4, 6, 2 so với 2 & 0,50 / 0,19 / 1 & mô tả \\
H4 & Xác nhận $\le 10\%$, chặn nhầm $\le 15\%$, token $\le +20\%$, độ trễ $\le +15\%$ & tối đa 0,6\% / 3,4\% /
  +0,8\% / +0,6\% & -- & mô tả, trong ngưỡng \\
\bottomrule
\end{tabularx}
\end{table}


\paragraph{Đánh đổi giữa an ninh, tiện ích và chi phí.} Trên mặt phẳng gồm số họ bị hại và $\Delta U$ ở
\ifpareto Hình~\ref{fig:pareto}\else Bảng~\ref{tab:main}\fi, các bản của \ours{} nằm ở vùng ít họ bị hại nhất, và không cấu
hình nào khác vừa ít họ bị hại hơn vừa giữ nhiều tiện ích hơn. Cái giá về tiện ích tăng dần khi luật chặt hơn, từ 0,8
điểm với \strRules{} lên 2,6 điểm với \strSand{} và 3,1 điểm với \ours. Chuyển từ \strSand{} sang \ours{} giảm thêm 4 họ
bị hại với cái giá là 5 ca lành tính, tương ứng 0,5 điểm. Vì \strSand{} giữ nhiều tiện ích hơn và bản bỏ \pThree{} giữ
hơn \ours{} một ca, không thể khẳng định \ours{} trội Pareto so với hai cấu hình này. Phòng thủ phát hiện gần như không
giảm được tác hại, trong khi phòng thủ bằng lời nhắc giảm được tác hại nhưng mất 2,1 đến 3,1 điểm tiện ích. Ngay cả chặn
mọi sink vẫn để 14 họ bị hại qua câu trả lời cuối và chỉ giữ 11,5\% tiện ích.

\ifpareto% Hình 2, pgfplots (Overleaf tự vẽ, không render trên máy). KHÔNG chạy regex đổi dấu thập phân trên file này: tọa độ
% dùng dấu chấm, nhãn trục dùng dấu phẩy nhờ /pgf/number format/use comma.
% Số liệu: facts.md mục Test-4 (họ bị hại /48, tiện ích) và mục "Bổ sung cho bài v3" (tiện ích dạng số ca /960).
% Delta U = (k - 411)/960*100 điểm, k = số ca lành tính thành công:
%   NoDef 411 -> 0 | MELON* 411 -> 0 | PromptGuard 410 -> -0.1042 | Progent* 409 -> -0.2083 | Sandwich 391 -> -2.0833
%   Spotlight 381 -> -3.125 | PhraseRules 412 -> +0.1042 | BG-Rules 403 -> -0.8333 | BG-Base 386 -> -2.6042
%   BrokerGuard 381 -> -3.125 | -C1 384 -> -2.8125 | -C2 386 -> -2.6042 | -C3 382 -> -3.0208 | chặn mọi sink 110 -> -31.35
% Màu theo thực thể như Hình 1, 3: xám = phòng thủ khác, xanh = dòng phát triển, cam = BrokerGuard (rỗng = bỏ một thành phần).
% Hình dạng điểm khác nhau nên đọc được khi in đen trắng. Tổng rộng khoảng 10,4 cm.
\begin{figure}[!htbp]
\centering
\begin{tikzpicture}
\begin{axis}[
  scale only axis, width=9.4cm, height=4.4cm,
  xmin=-1.5, xmax=42, ymin=-3.6, ymax=0.6,
  xtick={0,10,20,30,40}, ytick={-3,-2,-1,0},
  xlabel={Số họ bị hại trên \tThree{} (/48, gộp ba mô hình nền)},
  ylabel={$\Delta U$ so với \cNone{} (điểm \%)},
  axis x line*=bottom, axis y line*=left,
  tick label style={font=\scriptsize}, label style={font=\scriptsize},
  xmajorgrids, ymajorgrids, grid style={line width=0.2pt, draw=black!12},
  /pgf/number format/use comma,
  legend style={font=\tiny, draw=black!20, fill=white, fill opacity=0.9, text opacity=1,
    at={(0.99,0.56)}, anchor=north east, cells={anchor=west}},
  clip=false,
]
\addplot[only marks, mark=*, mark size=1.8pt, mark options={draw=black!70, fill=black!25}]
  coordinates {(40,0) (40,-0.1042) (38,-3.125) (22,-2.0833) (25,-0.2083)};
\addlegendentry{Phòng thủ khác}
\addplot[only marks, mark=square*, mark size=1.7pt, mark options={draw=bgblue, fill=bgblue}]
  coordinates {(34,0.1042) (15,-0.8333) (6,-2.6042)};
\addlegendentry{Dòng phát triển}
\addplot[only marks, mark=diamond, mark size=2.6pt, mark options={draw=bgorange, fill=white, line width=0.7pt}]
  coordinates {(4,-2.8125) (6,-2.6042) (2,-3.0208)};
\addlegendentry{\ours{} bỏ một thành phần}
\addplot[only marks, mark=diamond*, mark size=3pt, mark options={draw=bgorange!75!black, fill=bgorange}]
  coordinates {(2,-3.125)};
\addlegendentry{\ours}
% nhãn từng điểm
\node[font=\tiny, anchor=north east] at (axis cs:41.2,-0.22) {\cNone, \cMel, \cPG};
\node[font=\tiny, anchor=south] at (axis cs:34,0.2) {\cPhrase};
\node[font=\tiny, anchor=south] at (axis cs:25,-0.1) {\cProg};
\node[font=\tiny, anchor=south] at (axis cs:22,-1.97) {\cSand};
\node[font=\tiny, anchor=south] at (axis cs:38,-3.0) {\cSpot};
\node[font=\tiny, anchor=west] at (axis cs:15.8,-0.8333) {\cRules};
\node[font=\tiny, anchor=west] at (axis cs:6.9,-2.55) {\base, bỏ \cTwo};
\node[font=\tiny, anchor=west] at (axis cs:4.9,-2.86) {bỏ \cOne};
\node[font=\tiny, anchor=north west] at (axis cs:2.6,-3.17) {\textbf{\ours}, bỏ \cThree};
% hướng tốt hơn và tham chiếu ngoài trục (mũi tên đặt đúng x = 14 họ)
\draw[-{Stealth[length=3pt]}, draw=black!45, line width=0.5pt] (axis cs:11,-1.35) -- (axis cs:6.5,-0.55)
  node[anchor=south, font=\tiny, text=black!60] {tốt hơn};
\draw[-{Stealth[length=2.5pt]}, draw=black!45, line width=0.5pt] (axis cs:14,-3.0) -- (axis cs:14,-3.55);
\node[font=\tiny, text=black!60, anchor=west] at (axis cs:14.6,-3.3)
  {Chặn mọi sink: 14 họ, $\Delta U = -31{,}4$ (ngoài trục)};
\end{axis}
\end{tikzpicture}
\caption{Đánh đổi giữa an ninh và tiện ích trên \tThree. Trục ngang là số họ bị hại, trục dọc là $\Delta U$ so với
\cNone{} tính trên 960 ca lành tính. Góc trên bên trái là tốt nhất. Dòng phát triển gồm \cPhrase, \cRules{} và \base.
Điểm của \base{} trùng với \ours{} bỏ \cTwo, còn \ours{} bỏ \cThree{} giữ hơn \ours{} một ca lành tính. Tham chiếu chặn
mọi sink được ước lượng từ quỹ đạo \cNone{} và nằm ngoài trục.}\label{fig:pareto}
\end{figure}
\fi

Chi phí của ba thành phần mới gần như không đáng kể. Trên từng mô hình nền, số token trung bình mỗi ca của \ours{} chênh
không quá 0,8\% và thời gian trung bình mỗi ca chênh từ $-1{,}4\%$ đến $+0{,}6\%$ so với \strSand. Tỉ lệ cần xác nhận
trên ca lành tính tối đa là 0,6\%, tỉ lệ chặn nhầm từ 1,9\% đến 3,4\% so với 1,6\% đến 2,8\% của \strSand, đều dưới
ngưỡng của H4. Cổng chỉ mất trung bình 0,26~ms mỗi lời gọi, trong khi thời gian mô hình khoảng 14~s mỗi ca.

\paragraph{Phân tích theo cơ chế.} Bảng~\ref{tab:mechres} cho số họ bị hại theo cơ chế. \cSand{} và \cProg{} giảm mạnh
hai cơ chế chuyển hướng nhưng vẫn để lọt cả 8 họ rò rỉ qua sink, và \cProg{} cũng không giảm được rò rỉ qua câu trả lời
hay sửa tham số. Ngược lại, mỗi thành phần của \ours{} chặn đúng loại tấn công mà nó được thiết kế để chặn: bỏ \pTwo{}
làm 4 họ rò rỉ lọt trở lại, bỏ \pOne{} làm thêm 1 họ chuyển hướng và 1 họ sửa tham số bị hại, trong khi bỏ \pThree{} chỉ
thêm 1 ca. Ở mức ca, cả 4 họ chỉ bị hại dưới \strSand{} đều là vi phạm tại sink và đều bị cổng của \ours{} chặn trước khi
thực thi, 2 họ bởi \pTwo{} và 2 họ bởi \pOne. Khi bỏ đúng thành phần đó, ca tương ứng lại vượt qua. Xét 720 ca tấn công,
hai cấu hình chỉ khác nhau ở vi phạm tại sink, với 4 ca dưới \strSand{} và không ca nào dưới \ours. Ở cả hai, rò rỉ qua câu
trả lời đã về 0 và sửa tham số còn 5 ca, thuộc hai họ sửa giá trị số hoặc đơn vị mà cả hai cổng đều cho qua.

% Bảng 4 (bản 15 trang: gộp hình theo cơ chế và bảng thích nghi của bản dài). Nguồn: facts.md mục Test-4, bảng "Họ bị
% hại theo cơ chế (gộp 3 backbone, mỗi cơ chế /8)": A0 6 5 8 7 8 6 | P2 2 0 8 4 5 3 | PRG 1 0 8 7 8 1 | v1 1 0 1 1 3 0 |
% v3 0 0 0 0 2 0 | v3-C1 1 0 0 0 3 0 | v3-C2 0 0 2 2 2 0 | v3-C3 0 0 0 0 2 0 (thứ tự ER CR SL AL PT US = M1..M6);
% dòng "24 ca thích nghi" (ca hại Q/G/L) và mục "Vòng thích nghi nhắm v3 trên Dev-val" (họ bị hại /42 gộp ba mô hình,
% thích nghi · gốc). Tổng mỗi dòng theo cơ chế = cột Họ của Bảng 3. Bọc adjustbox max width.
\begin{table}[!htbp]
\centering
\caption{Số họ bị hại theo cơ chế trên \test{} (/8, gộp ba mô hình nền) và dưới hai bộ tấn công thích nghi. Bộ 1:
số ca gây hại trên 24 ca với Qwen3.5 / Gemma-4 / Llama-3.1. Bộ 2: số họ bị hại trên 42 họ của phần kiểm định, với
tấn công thích nghi $\cdot$ với tấn công gốc. Tên cơ chế theo Bảng~\ref{tab:mechanisms}.}
\label{tab:mechres}
\footnotesize
\setlength{\tabcolsep}{3pt}
\begin{adjustbox}{max width=\textwidth}
\begin{tabular*}{\textwidth}{@{\extracolsep{\fill}}lccccccrr@{}}
\toprule
 & \multicolumn{6}{c}{Họ bị hại theo cơ chế (/8)} & Bộ 1 (/24) & Bộ 2 (/42) \\
\cmidrule(lr){2-7}\cmidrule(l){8-9}
Cấu hình & \mER & \mCR & \mSL & \mAL & \mPT & \mUS & Q / G / L & thích nghi $\cdot$ gốc \\
\midrule
\NoDef & 6 & 5 & 8 & 7 & 8 & 6 & 8 / 6 / 7 & 25 $\cdot$ 26 \\
\cSand & 2 & 0 & 8 & 4 & 5 & 3 & 0 / 0 / 1 & -- \\
\cProg & 1 & 0 & 8 & 7 & 8 & 1 & 8 / 6 / 6 & -- \\
\StrSand & 1 & 0 & 1 & 1 & 3 & 0 & 0 / 0 / 1 & 2 $\cdot$ 2 \\
\textbf{\ours} & 0 & 0 & 0 & 0 & 2 & 0 & \textbf{0 / 0 / 0} & 1 $\cdot$ 1 \\
\quad bỏ \cOne & 1 & 0 & 0 & 0 & 3 & 0 & 0 / 0 / 0 & -- \\
\quad bỏ \cTwo & 0 & 0 & 2 & 2 & 2 & 0 & 0 / 0 / 1 & -- \\
\quad bỏ \cThree & 0 & 0 & 0 & 0 & 2 & 0 & 0 / 0 / 0 & -- \\
\bottomrule
\end{tabular*}
\end{adjustbox}
\end{table}


\paragraph{Tấn công thích nghi.} Ở bộ thứ nhất, \ours{} không bị hại ở ca nào, nhưng kết quả này thiên vị \ours{} vì bộ
này nhắm đúng vào các điểm yếu đã dùng làm đầu vào cho thiết kế. Khi không phòng thủ, bộ thứ hai gây hại ở 25 trên 42
họ, tương đương 26 họ của tấn công gốc, nên đây là các tấn công có hiệu lực. Với \strSand{} và \ours, không họ nào vượt
qua bằng tấn công thích nghi mà không vượt qua bằng tấn công gốc. Họ duy nhất còn vượt qua \ours{} trên Qwen3.5 thuộc cơ
chế sửa tham số, cùng loại với hai họ còn lại trên \test. Ở bộ này, tới 18 trong 42 ca trên Qwen3.5 dưới \strSand{} và
\ours{} dừng vì hết tám bước, nên kết quả chỉ có giá trị trong giới hạn số bước đó.

\paragraph{Tổng quát hóa.} Với Gemini 3 Flash trên \test, số họ bị hại là 26 khi không phòng thủ, 14 với \cProg{} và 0
với \cSand, \strSand{} và \ours, nên cột tham khảo này không phân biệt được hai cấu hình sau cùng về an ninh. \ours{} giữ
tiện ích \bPlain{} ở mức 201/240, so với 205/240 của \strSand, nhưng chặn nhầm 8 trong 320 ca lành tính, so với 4 ca của
\strSand.

Bảng~\ref{tab:agentdojo} cho kết quả trên AgentDojo. \ours{} cho cùng kết quả với \strSand{} ở Banking trong cả hai ngôn
ngữ. Ở Slack tiếng Anh, số tấn công thành công giảm từ 6 xuống 5: nhờ đòi khớp nguyên token, \pOne{} chặn được ca có tên
người nhận chỉ xuất hiện như một chuỗi con trong $u$. Các tấn công còn thành công gồm đổi mật khẩu qua một sink không có
đích, với 3 ca tiếng Anh và 1 ca tiếng Việt ở Banking, mở một URL bằng công cụ đọc ở 4 ca, và 1 ca khác ở Slack tiếng
Anh. Hai loại đầu nằm ngoài tập sink $S$ của mô hình hình thức. Đổi lại, tiện ích ở Slack giảm thêm, dưới tấn công từ 15
xuống 11 ở tiếng Anh và khi không có tấn công từ 5 xuống 3 trên 21 nhiệm vụ ở cả hai ngôn ngữ, vì nhiều nhiệm vụ Slack lấy
đích từ một trang web hay tin nhắn mà người dùng không viết.

% Bảng 5. facts.md mục "AgentDojo: bảng đầy đủ mọi lượt" (NoDef, repeat_user_prompt = Sandwich, KT) và "AgentDojo:
% BrokerGuard v3". Ô = utility dưới tấn công / tấn công thành công. Bản 15 trang: phần Gemini của bảng tổng quát hóa cũ
% (facts.md "Cột tham khảo Gemini 3 Flash trên Test-4") chuyển thành lời văn ở Mục 7.4. Bọc adjustbox max width.
\begin{table}[!htbp]
\centering
\caption{AgentDojo với Llama-3.1-8B và khuôn \code{important\_instructions}. Ô là tiện ích dưới tấn công / số tấn công
thành công (Banking /144, Slack /105). \cSand{} là \code{repeat\_user\_prompt} của AgentDojo. Cột cuối là tiện ích khi
không tấn công (Banking /16, Slack /21) khi không phòng thủ, với \strSand{} và với \ours.}\label{tab:agentdojo}
\footnotesize
\setlength{\tabcolsep}{3pt}
\begin{adjustbox}{max width=\textwidth}
\begin{tabular*}{\textwidth}{@{\extracolsep{\fill}}llccccc@{}}
\toprule
Bộ & Ngôn ngữ & \makecell{Không\\phòng thủ} & \cSand & \makecell{Luật chuỗi\\+ Sandwich} & \ours &
  \makecell{Không\\tấn công} \\
\midrule
Banking & tiếng Anh & 58/17 & 62/10 & 72/3 & 72/3 & 7 / 8 / 8 \\
Banking & tiếng Việt & 44/6 & 45/6 & 54/1 & 54/1 & 6 / 5 / 5 \\
Slack & tiếng Anh & 28/16 & 35/10 & 15/6 & 11/5 & 7 / 5 / 3 \\
Slack & tiếng Việt & 26/0 & 19/0 & 13/0 & 12/0 & 7 / 5 / 3 \\
\bottomrule
\end{tabular*}
\end{adjustbox}
\end{table}


% Số: facts.md mục Test-4 ("Hạn chế phải nêu (v3)": lực 26%, ~120 họ, đọc trái phép v1 20 / v3 18 ca, tiện ích tuyệt
% đối thấp; "40 ca vượt v1" trong mẫu 38 -> 2; M3 Test-II 9 cặp so với Test-4 1), mục "Bổ sung cho bài v3" (phát lại
% oracle Dev-val), mục baseline bổ sung (Prompt Guard 0 lần chặn), mục AgentDojo v3.
\section{Thảo luận}\label{sec:discussion}

\paragraph{Vì sao giả thuyết chính không đạt.} H1 được thiết kế với giả định \base{} để lọt khoảng 14 trong 48 họ,
như trên \tTwo. Trên \tThree, \base{} chỉ để lọt 6 họ. Khác biệt nằm chủ yếu ở rò rỉ qua sink: trên \tTwo{} cơ chế này
có 9 cặp họ và mô hình bị hại dưới \base, trên \tThree{} chỉ còn 1. Với nền thấp như vậy, ngay cả khi \ours{} chặn được
mọi họ mà \base{} để lọt, số cặp lệch vẫn rất nhỏ. Lực của kiểm định ở mức nền này chỉ khoảng 26\%, và cần khoảng 120
họ để đạt lực 80\%. Kết quả trên \tThree{} vì vậy là một hiệu ứng sàn chứ không phải bằng chứng rằng hai cấu hình như
nhau. Các nguồn khác cho cùng chiều. Trên \dev, số cặp họ và mô hình bị hại giảm từ 36 xuống 4. Chấm lại 40 ca từng lọt
qua \base{} trên hai tập niêm phong trước cho 38 ca bị hại dưới \base{} và 2 dưới \ours, cả hai đều là sửa tham số.
Tuy vậy, các số này đến từ dữ liệu đã dùng khi phát triển \ours, nên chỉ có giá trị chẩn đoán và không thay được kiểm
định đăng ký trước. Bài báo không thêm mô hình hay chạy lại để tìm ý nghĩa thống kê sau khi đã biết kết quả.

\paragraph{Mỗi thành phần đóng góp gì.} Bằng chứng mạnh nhất cho \cOne{} và \cTwo{} là bằng chứng mức ca. Bốn họ mà
\base{} để lọt đều bị cổng của \ours{} chặn trước khi thực thi, và mỗi họ lọt trở lại khi bỏ đúng thành phần đã chặn
nó. \cTwo{} xử lý đúng chỗ hở lớn nhất của \base, là giá trị mật nằm trong nguồn gắn nhãn công khai. Dưới \ours, không
giá trị nào lộ ở câu trả lời cuối. Tác dụng của \cOne{} khi chạy thật nhỏ hơn khi chấm offline. Trên các lời gọi sink
của quỹ đạo tham chiếu thuộc phần kiểm định của \dev, \cOne{} chặn đủ 35 lời gọi tác dụng phụ ngoài yêu cầu, so với 15
của \rDest. Nhưng các mô hình hiếm khi tự đề xuất những đích này, nên trên \tThree{} \cOne{} chỉ tạo khác biệt ở 2 họ.
\cThree{} chưa có tác dụng đo được trên \tThree. Hai họ còn lọt là sửa giá trị số hoặc đơn vị. Có hai cách giải thích, đều chưa được kiểm từng ca.
Hoặc giá trị bị sửa không nằm trong tập literal được ghim, tức (A5) bị vi phạm. Hoặc literal vẫn có mặt cạnh một giá
trị khác, trường hợp Mệnh đề~\ref{prop:int} không loại trừ. Cả hai đều cho thấy \cThree{} cần ngữ nghĩa của trường
thay vì chỉ kiểm sự có mặt của literal.

\paragraph{Chi phí tiện ích.} \ours{} không làm tiện ích thấp hơn \base{} quá 0,6 điểm ở nhánh nào, nhưng cả hai đều
thấp hơn \cNone{} khoảng 3 điểm. Một phần nguyên nhân là chặn nhầm. \ours{} chặn nhầm 24 trong 960 ca lành tính của
\tThree, so với 19 của \base. Trên \dev, lý do chặn nhầm thường gặp là các luật giữ từ \base{} (\rFlow{} theo nguồn và
\rSig). Riêng với Llama-3.1, \cThree{} đứng đầu, vì mô hình này hay viết lại giá trị người dùng đã nêu khi soạn nội
dung gửi. Tiện ích tuyệt đối thấp ở mọi cấu hình, kể cả \cNone{} (42,8\%), do giới hạn của các mô hình nhỏ như hết bước
hay đầu ra bị cắt, không do phòng thủ. Trên AgentDojo, chi phí tập trung ở các nhiệm vụ lấy đích từ dữ liệu
công cụ. Trong benchmark, ``cần xác nhận'' được thi hành như từ chối. Một triển khai thật có thể hỏi người dùng ở bước
này, và đây là cách hợp lý để lấy lại tiện ích mà không cấp quyền cho nội dung không đáng tin.

\paragraph{Ngôn ngữ.} Luật \rDest{} của \base{} so khớp chuỗi nên không phụ thuộc ngôn ngữ. \cOne{} thì có, vì cực tính
dựa trên danh sách từ phủ định, loại trừ và tương phản cho tiếng Việt có dấu, không dấu và tiếng Anh. Một ngôn ngữ hay
phương ngữ mới cần bổ sung danh sách này, và một phủ định ngoài danh sách sẽ làm cổng cho qua một đích lẽ ra bị cấm.
Thiết kế giảm rủi ro này bằng cách trả về \emph{cần xác nhận} khi phủ định không rõ phạm vi, nhưng không loại bỏ được
nó. Với phòng thủ phát hiện, vấn đề ngôn ngữ còn rõ hơn: \cPG{} không gắn cờ một quan sát nào của \tThree, tức bộ phân
loại này không chuyển được sang dữ liệu tiếng Việt của benchmark.

\paragraph{Mối đe dọa với tính hợp lệ.} Về cấu trúc đo, vị từ gây hại và bộ chấm được đóng băng trước khi chạy, nhưng
dữ liệu là tổng hợp, công cụ là bản giả và mỗi nhiệm vụ chỉ có một lượt. Về tính hợp lệ nội tại, \tThree{} được sinh
từ đặc tả bằng mã và được chính nhóm phát triển phòng thủ rà soát. Tập được niêm phong trước khi có mã của \ours, nhưng
không loại trừ được việc người soạn vô tình chọn các ca hợp với thiết kế. Bộ thích nghi thứ hai do thành viên của nhóm
soạn, chỉ có một vòng và chạy trên tập phát triển. Năm baseline bổ sung được chạy sau khi đã biết kết quả chính. Về
tính hợp lệ bên ngoài, các mô hình mở chỉ có 8 đến 12 tỉ tham số, Gemini 3 Flash chạy qua proxy và không ghim được
phiên bản, còn AgentDojo chỉ chạy với một mô hình.

\paragraph{Hạn chế.} Các mệnh đề được chứng minh trên đặc tả, chưa được kiểm hình thức trên mã. Kiểm thử dựa trên thuộc
tính với một văn phạm sinh yêu cầu và chuỗi lời gọi là bước tiếp theo tự nhiên. Đọc trái phép nằm ngoài phạm vi: dưới
\ours{} vẫn có 18 trong 720 ca tấn công có lời gọi đọc trái phép, gần như toàn bộ trên Llama-3.1 (\base: 20). \cTwo{}
không theo dõi giá trị qua các phép biến đổi ngoài (A4), như mã hóa hay chia nhỏ qua nhiều lời gọi. Bộ chấm biến đổi
không tìm thấy ca nào như vậy trên dữ liệu hiện có, nên giả định này chưa thực sự bị thử thách. Cuối cùng, bộ dò của
\cTwo{} được thiết kế để không phụ thuộc vào dạng giá trị mật của bộ sinh dữ liệu, nhưng độ phủ trên giá trị mật thật
vẫn phải đo trên dữ liệu thật.

% Số: như tóm tắt (facts.md mục Test-4, AgentDojo v3). Hướng tiếp theo nối với các hạn chế ở Mục 8.
\section{Kết luận}\label{sec:conclusion}

Bài báo trình bày \ours, một cơ chế cưỡng chế tại Tool Broker gồm chính sách ý định có cực tính, nhãn độ nhạy cảm mức
giá trị và ghim tham số, cùng một mô hình hình thức với ba thuộc tính được chứng minh dưới năm giả định. Trên một tập
kiểm thử tiếng Việt niêm phong trước khi có mã, \ours{} giảm số họ bị hại từ 40 khi không phòng thủ xuống 2 trên 48,
với tiện ích giảm 3,1 điểm, còn token và độ trễ mỗi ca tăng không quá 1\% so với \strSand. So với cấu hình này (6 họ),
mức giảm chưa đạt ý nghĩa thống kê do hiệu ứng sàn, nhưng mỗi họ mà luật so khớp chuỗi để lọt đều bị chặn bởi đúng
thành phần được thiết kế cho nó. Hướng tiếp theo gồm một tập kiểm thử niêm phong khoảng 120 họ để có đủ lực thống kê,
kiểm hình thức các mệnh đề trên mã, ngữ nghĩa trường cho \pThree, và kênh xác nhận với người dùng cho đích lấy từ dữ
liệu công cụ.


% \begin{credits}
% \subsubsection{\ackname} ...
% \subsubsection{\discintname} Các tác giả không có xung đột lợi ích.
% \end{credits}

\bibliographystyle{splncs04}
\bibliography{refs}

\end{document}
